% Source: https://artofproblemsolving.com/community/c3h1615036_digital_roots
% Stable identity: aops-topic-1615036-q01
A number and its digit sum leave the same remainder on division by $9$, since $10\equiv1\pmod 9$ gives $\overline{d_k\dots d_1d_0}=\sum d_i10^i\equiv\sum d_i\pmod 9$. So every step of the process keeps the remainder mod $9$. The final single digit is between $1$ and $9$, so the digital root is $1$ exactly when the number is $\equiv1\pmod 9$.

The three-digit numbers that are $\equiv1\pmod9$ are $100,109,118,\dots,991$, an arithmetic sequence with difference $9$ ($100=9\cdot11+1$ and $991=9\cdot110+1$). Their number is
\[
\dfrac{991-100}{9}+1=99+1=\boxed{100}.
\]
This confirms the count $100$ obtained in the post by listing.
